\documentclass[11pt]{article}
\usepackage[review]{acl}
\usepackage{times}
\usepackage{latexsym}
\usepackage[T1]{fontenc}
\usepackage[utf8]{inputenc}
\usepackage{microtype}
\usepackage{inconsolata}
\usepackage{graphicx}
\usepackage{float}
\usepackage[table]{xcolor}
\usepackage{array,booktabs,amsmath,amssymb,multirow,url}

\title{ColorRef: Interactive Reference Games for Probing the Behavioral Geometry of Color Grounding}
\author{Pablo Loyola\textsuperscript{1}, Andres Hoyos-Idobro\textsuperscript{2}, Edison Marrese-Taylor\textsuperscript{3} \\
\textsuperscript{1} Rakuten Institute of Technology; Rakuten Group, Inc; Tokyo, Japan\\
\textsuperscript{2} Rakuten Institute of Technology; Rakuten Group, Inc; Paris, France\\
\textsuperscript{3} Graduate School of Engineering, The University of Tokyo\\
\{pablo.a.loyola,andres.hoyosidrobo\}@rakuten.com, emarrese@weblab.t.u-tokyo.ac.jp}

\begin{document}
\maketitle

\begin{abstract}
We introduce \textsc{ColorRef}, an interactive reference-game framework for
studying natural-language control of color predictions. Using 518{,}830
English color descriptions, we separate initial prediction, output encoding,
direction following, correction magnitude, and teacher feedback.
Controlled feedback reduces Qwen3-14B's mean target error from 42.33 to
22.15 $\Delta E_{76}$ on 1{,}000 new regime-balanced descriptions.
Matched-start controls show sensitivity to output representation and explicit
axis instructions. Graded corrections exhibit ordered but direction-dependent
step sizes; correct directions alone do not prevent overshoot.
A frozen wording policy calibrated on 32 separate colors reduces one-step
held-out target error from 12.36 to 6.86 relative to a fixed graded-wording
rule across 128 starting colors. A smaller sequential pilot also tests
transfer and generation cost.
Oracle-assisted teacher controls further distinguish preserving supplied
directions from producing useful corrections: few-shot restatement can be
faithful without yielding the best receiver accuracy.
These results characterize prompt-conditioned control of recorded color
targets, not uniquely correct color meanings or human perceptual calibration.
\end{abstract}

\section{Introduction}
\label{sec:introduction}

Color descriptions connect language to a continuous perceptual domain.
Some expressions directly name a region, such as \textit{dark green};
others invoke objects, materials, or associations, such as \textit{ocean blue}
or \textit{abandoned hope}. A model's initial color prediction is informative,
but grounded communication also involves revision. A user may request
\textit{more blue}, then \textit{a little lighter}: the system must identify
both a direction and an appropriate magnitude relative to its current state.

Most language-to-color evaluations test one-shot prediction or reference
among alternatives \citep{mcmahan2015bayesian,monroe2017colors}.
Representation studies ask whether model embeddings align with perceptual
color geometry \citep{abdou-etal-2021-language,loyola-etal-2023-perceptual}.
We ask a complementary behavioral question: how reliably can language steer
model outputs through color space? This requires separating initial prediction
from output encoding, direction interpretation, step size, and feedback
production. A geometrically valid direction need not reduce error if the
receiver moves too far; a faithful teacher utterance need not induce a useful
update.

\textsc{ColorRef} makes these distinctions measurable. A guesser predicts a
color from a description; a teacher sees the hidden target and provides
corrections. We compare deterministic oracle policies and model-produced
feedback, varying constraints per message, wording, and interaction length.
Four semantic regimes distinguish direct color terms from prototype-mediated,
compound-associative, and abstract-idiosyncratic descriptions.

We report $\Delta E_{76}$, the Euclidean distance between predicted and
target colors in CIELAB space \citep{colorimetry1986official}.
Lower values indicate closer agreement with the recorded target, not that
alternative colors are semantically invalid. This distinction matters for
creative or ambiguous names. We use the metric as a reproducible distance,
not a strict just-noticeable-difference threshold.

Our fresh evaluation shows that controlled feedback reduces mean disagreement
with recorded targets by approximately 48\%. We complement this result with
controls that hold starting
colors fixed, probe graded instructions independently of color names, and
evaluate a calibration-derived wording policy on unseen starting colors.
The latter demonstrates controlled target localization relevant to iterative
creative interfaces; image editing and human interaction remain untested.

Finally, we compare zero- and four-shot oracle-assisted teachers under natural
and restricted wording instructions, then pass their unchanged messages to
the same receiver. The results replace a broad following-versus-generating
capability claim with a more specific finding: feedback preservation, output
format, and downstream usefulness are distinct, prompt-sensitive outcomes.

\section{Related Work}
\label{sec:related_work}

\paragraph{Color semantics and representation geometry.}
Color language has been studied through cross-linguistic regularities,
perceptual geometry, and communicative efficiency
\citep{berlin1991basic,regier2007color,zaslavsky2018efficient}.
Probabilistic and neural approaches map descriptions to perceptual
observations \citep{mcmahan2015bayesian,monroe2017colors}.
\citet{abdou-etal-2021-language} compare color-term representations with
CIELAB geometry; \citet{loyola-etal-2023-perceptual} study how alignment
changes with subjectivity and abstraction. We complement representation
probing with controlled output trajectories rather than infer internal
representations from behavioral success.

\paragraph{Comparatives and interactive grounding.}
\citet{winn-muresan-2018-lighter} model comparative color terms as directions
in color space. Reference games and pragmatic models treat communication as
speaker--listener coordination
\citep{lewis1969convention,frank2012predicting,goodman2016pragmatic}.
Grounded clarification distinguishes perceptual and collaborative grounding
\citep{benotti2021recipe}, while interactive evaluation studies common-ground
formation \citep{imai2025measuring}. Our context is the current continuous
state: we examine both directional corrections and the operational magnitude
of graded modifiers.

\paragraph{Feedback and revision.}
Feedback supports instruction tuning, AI critique, and iterative refinement
\citep{ouyang2022training,bai2022constitutional,saunders2022selfcritiquing,
madaan2023selfrefine,shinn2023reflexion}. However, feedback and its reception
can introduce separate errors. We use oracle-assisted restatement and matched
receiving revisions to distinguish supplied-direction preservation from
usefulness, without equating linguistic audit scores with task success.

\section{Framework}
\label{sec:colorref}

\subsection{Descriptions and assigned targets}
\label{sec:dataset}
We construct the corpus from ColorNames-derived English descriptions and
sRGB hex codes. After preprocessing and validity filtering, 518{,}830 rows
remain. Each example $(x_i,h_i,c_i,r_i)$ contains a description, recorded
HEX target, its CIELAB coordinate, and a semantic-regime label.
We retain description-validity scores $\geq 0.75$, including 43{,}187 rows
at the boundary. The validity filter asks whether a string plausibly functions as a color
description; it does not measure agreement about its paired target.
Construction and filtering details appear in Appendix~\ref{sec:appendix-dataset}.

The four regimes are \emph{explicit grounded}, with direct color terms;
\emph{prototype mediated}, referring through objects or materials;
\emph{compound associative}, combining a color anchor with other content;
and \emph{abstract idiosyncratic}, using metaphorical or affective expressions.
They are coarse parser-based categories, not exhaustive semantic analyses.
Examples and counts appear in Appendix~\ref{app:corpus-examples}.

\subsection{Interactive states and evaluation}
\label{sec:game}
In the original protocol, the guesser initially observes only $x_i$ and
outputs HEX. The teacher sees the target and current guess, emits feedback,
and the guesser receives the description, previous guess, and latest message.
The reference game permits up to three revisions, with early stopping at
$\Delta E_{76}<5$. Controlled studies below either supply the start or
replace description prediction with an isolated adjustment task.

For direct-LAB outputs, retain both native coordinates $z_t$ and the displayed
state $\hat c_t=P(z_t)$, obtained by conversion to clipped, rounded uint8
sRGB and back to LAB. Feedback and subsequent prompts use the displayed state
in the matched controls. Native and displayed behavior can differ:
rectangular LAB encoding bounds are not the sRGB gamut.

\subsection{Metrics}
\label{sec:metrics}
The primary error is $e_t=\|c-\hat c_t\|_2=\Delta E_{76}$.
We report initial and endpoint error, gain $e_0-e_T$, and mean per-example
relative improvement $(e_0-e_T)/(e_0+\epsilon)$; the latter is not the
percentage reduction of the cohort mean. Alignment compares the displayed
update $u_t=\hat c_{t+1}-\hat c_t$ with ideal correction $v_t=c-\hat c_t$:
\begin{equation}
A_t=\frac{v_t^\top u_t}{\|v_t\|_2\|u_t\|_2}.
\end{equation}
Zero-length vectors are excluded from alignment, not from target error.
Constraint satisfaction checks signed movement along supplied oracle
directions, including HSV saturation where applicable.

For an isolated instruction with unit requested direction $d$, we measure
signed step $s=d^\top(z'-z)$ and off-axis drift
$\|(z'-z)-sd\|_2$, also reporting displayed counterparts.
Positive steps follow the requested direction; zero and wrong-sign steps
remain separate. Ordinal comparisons test the ladder
\emph{a little} $<$ \emph{somewhat} $<$ \emph{much}; bare wording is unranked.
Projection cost is $\|z'-P(z')\|_2$.
Direction compliance alone is insufficient: a nonzero update reduces
Euclidean error only when $\|u\|_2<2\|v\|_2\cos\theta$.

\section{Experimental Setup}
\label{sec:experimental_setup}

\subsection{Oracle feedback and fresh grounding evaluation}
Oracle policies compute LAB lightness, red--green, yellow--blue and HSV
saturation discrepancies, threshold candidates, and rank them by normalized
absolute error. They emit up to one, two, or three selected constraints.
The minimal and axis-$c1$ policies share a one-constraint setting; template
wording expresses the same selected directions through deterministic phrases.
These phrases do not introduce the graded modifiers studied separately below.
Examples appear in Appendix~\ref{app:bandwidth_wording_examples}.
An additional primitive-direction template variant was originally named
\emph{template-vocab}; the implementation does not establish a ceiling for
all broader teacher vocabulary.

The fresh Qwen3-14B grounding evaluation uses 1{,}000 descriptions, 250 per
regime, sampled with seed 113 after excluding all 400 original debug-subset
IDs. Every game uses HEX output and three fixed revisions with canonical
axis feedback supplying at most three constraints. The initial prediction
also provides its paired one-shot baseline. All 4{,}000 responses parse.
The sample is equally weighted across regimes, not weighted by corpus
frequency. Plans, selected IDs, source/subset hashes and prompts are frozen;
the model revision is pinned to \texttt{40c06982}. This protocol is a fresh
evaluation, rather than reconstruction of the earlier unrecoverable runs.

\subsection{Matched output-interface control}
\label{sec:setup_interface}
On 64 descriptions, 16 per regime, compare HEX, plain LAB, and LAB with an
explicit axis legend. Each receives the same displayed start from the parent
HEX predictions and the same first oracle correction, with at most three
constraints. Three fixed revisions yield 192 games and 576 generations;
later adaptive messages can differ. The first revision is the primary
matched endpoint, and the third is secondary. Fresh contexts show only the
latest feedback. Supplied starts cost no generation within this control.

\subsection{Graded wording and held-out control}
\label{sec:setup_magnitude}
An isolated probe uses 12 fixed colors, six directions, four wordings
(bare, a little, somewhat, much), and plain versus axis-legend LAB prompts.
Exact 5-/10-unit and no-change controls separate coordinate execution from
qualitative instructions: 888 responses in total, all in fresh contexts.

For policy calibration, 12 separately drawn sRGB colors produce 288
direction/wording responses. Fit direction-specific medians of displayed
signed movement, including nonpositive parsed movements; only positive
graded medians are selection candidates. For a hidden target, choose the
phrase whose median is closest to the current requested-axis residual.
The bare arm supplies direction only; a numeric control supplies the exact
coordinate increment and explicit holds on the other coordinates.
These arms share one requested axis, not equal semantic information.

The four-arm confirmation fits medians on 32 new colors (768 responses)
and evaluates 128 disjoint starting colors. An unfitted arm selects
\textit{a little}, \textit{somewhat}, or \textit{much} using fixed residual
cutpoints 9 and 18, without receiver medians. The primary comparison is
calibrated minus unfitted displayed target error. Input-only feasibility
checks retain 2{,}108 targets and exclude 196; previous study colors are
excluded before inference. Whole starting colors are resampled for intervals.

One-step evaluation uses 12 disjoint starts and 196 feasible targets.
Sequential transfer freezes the fitted policy and uses 12 further fresh
starts, 200 targets, target-known stopping at $\Delta E_{76}\leq5$, and a
five-revision cap. Targets are constructed at native distances 6, 12, or 24
along one of six LAB directions. Input-only feasibility checks exclude
20/216 and 16/216 targets, respectively; generated outputs are never filtered
by projection cost. The original axis remains fixed during transfer, but
the requested sign can reverse after overshoot. Appendix~\ref{app:new-controls}
details selection, projection, and cost accounting.

\subsection{Teacher restatement and feedback reception}
\label{sec:setup_teacher}
From the shared-start cohort, four input-selected demonstrations (one per
regime) leave 60 evaluation examples. Compare the archived natural-paraphrase
prompt with a restricted exact-direction prompt, each zero-shot and four-shot,
at maximum bandwidths one and three: 480 teacher messages.
All conditions see identical target/guess states and oracle directions;
the clause budget always covers the supplied count. Demonstration answers
are canonical oracle sentences, not teacher outputs selected by performance.
This tests oracle-assisted restatement, not unaided geometric generation.

A separate receiver study passes every original teacher message unchanged,
plus a canonical oracle arm, to the same HEX guesser from identical starts:
120 cases and 600 single revisions. No audit labels select or repair messages;
target coordinates and direction JSON are not appended to receiver prompts.
We keep frozen lexical audits, supplementary scope-aware interpretation,
and receiving target error separate.

\subsection{Models, uncertainty, and failures}
All new controls use \texttt{Qwen/Qwen3-14B} \citep{yang2025qwen3} through
Hugging Face Transformers \citep{wolf2020transformers}, bfloat16 on an
A100 40\,GB, temperature zero, and thinking disabled. Receiver budgets are
64 output tokens; teacher budgets are 80. Plans, prompts, model-revision
metadata, and atomic response checkpoints are preserved.
Strict parse failures remain failures, without repaired colors; backend
exceptions remain pending. All reported receiving/control evaluations here
parse fully; calibration parses 287/288 responses.

Original main intervals use 1{,}000 percentile bootstrap draws.
New paired intervals use 2{,}000 draws: whole examples within regimes for
interface/teacher reception, and whole starting colors for magnitude control,
retaining all associated cases and pooled weights. They condition on observed
cohorts, not generation randomness or unrestricted population generalization.
Subgroup analyses are exploratory and not multiplicity-adjusted.

\section{Results}
\label{sec:results}

\subsection{Interaction localizes recorded targets}
\label{sec:results_main}
Across all 1{,}000 fresh descriptions, mean initial error is 42.33
[40.70, 44.01] and final error is 22.15 [21.28, 22.99]. The paired
initial-minus-final gain is 20.19 [18.61, 21.79] $\Delta E_{76}$, reducing
mean error by approximately 47.7\%. This reduction of means is not mean
per-example relative improvement. Initial and final convergence counts
are 5 and 38 of 1{,}000 at the $\Delta E_{76}\leq5$ threshold.

\begin{table}[t]
\centering\small
\begin{tabular}{@{}lrrr@{}}
\toprule
Regime & Initial $\Delta E$ & Final $\Delta E$ & Gain \\
\midrule
Explicit & 36.41 & 22.02 & 14.39 \\
Prototype & 43.39 & 22.41 & 20.98 \\
Compound & 37.89 & 20.64 & 17.24 \\
Abstract & 51.64 & 23.51 & 28.12 \\
Overall & 42.33 & 22.15 & 20.19 \\
\bottomrule
\end{tabular}
\caption{Fresh Qwen3-14B grounding evaluation: 250 descriptions per regime,
canonical $c3$ feedback and three fixed revisions. Initial predictions are
the paired one-shot baseline. Gain is initial minus final error.}
\label{tab:main_results}
\end{table}

All regimes improve descriptively (Table~\ref{tab:main_results}). Abstract
descriptions have both greater initial disagreement and larger absolute
gain; these means do not establish a separate between-regime effect or
semantic implausibility of alternative colors. This evaluation measures
localization after additional corrections; it does not isolate sequential
adaptation from feedback information at a matched constraint budget.

\subsection{Output encoding affects correction from shared starts}
\label{sec:results_interface}
All 576 shared-start revisions parse; every interface is evaluated on all
64 examples (Table~\ref{tab:shared_interface}). Plain LAB exceeds HEX
first-revision error by 6.19 [0.85, 11.89] $\Delta E_{76}$.
Adding an axis legend reduces that error by 5.71, but its first-revision
interval spans zero. After three revisions, legend minus plain LAB is
$-12.11$ [$-18.04$, $-5.49$]. Thus, the encoding/instruction intervention
affects feedback following beyond differences in initial prediction.
Legend and HEX endpoints are close, but their difference 0.38
[$-4.85$, 6.29] does not establish equivalence. Later feedback is adaptive,
and the result is conditional on HEX-derived starts.

\begin{table}[t]
\centering\small
\begin{tabular}{@{}lrrr@{}}
\toprule
Interface & First $\Delta E$ & Third $\Delta E$ & C-sat \\
\midrule
HEX & 32.10 & 22.78 & 0.888 \\
Plain LAB & 38.29 & 35.28 & 0.823 \\
LAB + legend & 32.58 & 23.16 & 0.906 \\
\bottomrule
\end{tabular}
\caption{Shared starts and first corrections, $N=64$ per interface;
mean starting error 41.96. Displayed-state error is primary. Three revisions
are forced rather than stopped at convergence.}
\label{tab:shared_interface}
\end{table}

\subsection{Graded instructions specify unequal step sizes}
\label{sec:results_quantifiers}
All 888 isolated-probe responses parse. Explicit axis instructions increase
positive-direction qualitative updates from 228/288 to 288/288.
Strictly increasing ordinal comparisons rise from 172/216 to 210/216;
nondecreasing comparisons, which include ties, rise from 208/216 to 216/216.
These overlapping within-color comparisons concern 12 fixed anchors, not
216 independent colors. Numeric and no-change controls are exact except
one legend-condition 10-unit instruction executed as a 20-unit update.

The separate calibration split gives displayed median steps for
\emph{a little} ranging from 1.14 to 4.92 across directions,
\emph{somewhat} from 5.07 to 10.09, and \emph{much} from 28.72 to 75.84.
These are model/prompt-specific operational measurements, not universal
perceptual units. Large updates often leave the displayable gamut:
54/72 calibration responses for \emph{much} have projection error above one.
Native step variation and display projection must be distinguished; neither
the model's internal interpretation nor human perceptual nonuniformity is
identified by these outputs alone.

\subsection{Calibrated magnitude improves held-out localization}
\label{sec:results_magnitude}
The four-arm confirmation completes all 8{,}432 held-out responses.
There are 2{,}098 fully parsed quartets spanning all 128 starting colors;
parse failures number 3, 5, 0, and 2 for bare, unfitted, calibrated, and
numeric feedback. Calibrated wording reduces mean target error from 12.36
to 6.86 relative to the unfitted rule: paired difference $-5.50$
[$-5.79$, $-5.20$]. The available-pair contrast is similar ($-5.53$).
Unfitted versus bare error differs by $0.45$ [$-0.69$, $1.52$], so
graded wording alone does not establish improvement in this comparison.
Both graded arms follow the requested direction on approximately 99\%
of parsed responses. Direction compliance therefore does not ensure
accurate magnitude control. This tests the fitted policy against one fixed
rule; it does not establish that every unfitted rule is inferior.

Exploratory saved-output breakdowns localize the aggregate benefit. At
requested distances 6, 12, and 24, calibrated-minus-unfitted error is
$-0.71$, $0.00$, and $-18.35$, respectively. The largest reductions occur
for 24-unit red, yellow, and blue corrections; the 24-unit calibrated
mean error remains 15.31. Harmful updates fall from 318 to 58 on the
common cohort. However, 1{,}419 cases tie within 0.01, versus 556 wins
and 123 losses, and the median paired difference is zero. Calibration
chiefly avoids excessive large corrections in this task rather than
improving every adjustment. Subgroup intervals are exploratory and
the prompt-identity audit distinguishes equal errors from equal interventions.

\begin{table}[t]
\centering\small
\begin{tabular}{@{}llrrr@{}}
\toprule
Study & Arm & Final $\Delta E$ & $\leq5$ & Calls \\
\midrule
One step & Bare & 11.91 & -- & 2{,}108 \\
 & Unfitted & 12.36 & -- & 2{,}108 \\
 & Calibrated & 6.86 & -- & 2{,}108 \\
 & Numeric & 0.35 & -- & 2{,}108 \\
\midrule
Sequential & Bare & 6.46 & 148 & 516 \\
 & Calibrated & 2.98 & 192 & 338 \\
 & Numeric & 0.10 & 200 & 205 \\
\bottomrule
\end{tabular}
\caption{One-step confirmation: means on 2{,}098 common parsed quartets
from 128 starts; calls include failed outputs. Sequential pilot: 200 targets
from 12 further starts, stopping at error $\leq5$ with a five-revision cap.
Calls exclude calibration: 768 responses for confirmation and 288 for the
sequential pilot. Numeric feedback supplies additional precision.}
\label{tab:magnitude}
\end{table}

Transfer to 12 further starts completes all 600 games with 1{,}059 parsed
revisions. Calibrated minus bare endpoint error is $-3.48$
[$-4.20$, $-2.74$]; convergence rises from 148/200 to 192/200.
Evaluation calls fall from 516 to 338 (34.5\%), excluding calibration.
All twelve per-start mean contrasts favor the calibrated policy, but
individual targets do not uniformly improve. This sequential pilot has no
unfitted graded-selection arm; the larger four-arm confirmation tests that
contrast for one revision only.

Eight calibrated trajectories exhaust the cap: all concern 24-unit lighter
or greener targets and end with off-axis disagreement exceeding the
threshold. The fixed-axis policy does not explicitly restore display-induced
drift. Numeric feedback has the lowest endpoint error, but retains execution
tails: four one-step responses have native error above one (maximum 39),
and five sequential responses exceed one (maximum 44). The latter reach
convergence after subsequent oracle corrections, not autonomous self-repair.
Appendix~\ref{app:magnitude_diagnostics} preserves these limitations.

\subsection{Restatement fidelity and usefulness are distinct}
\label{sec:results_llm_teacher}
All 480 teacher messages and 600 receiver revisions complete.
A supplementary scope-aware audit resolves all supplied-direction sets in
both four-shot conditions as preserved; restricted zero-shot preserves
104/120. Natural zero-shot has 81 resolved matches, 17 mismatches, and
22 unresolved cases (Table~\ref{tab:teacher_receiver}).
These are finite-grammar diagnostics, not independent human judgments.
Literal keywords can invert the intended correction: \textit{your guess is
darker than the target} implies a lighter update. Harmless \textit{Feedback:}
prefixes can also lower strict format certification.

\begin{table}[t]
\centering\small
\begin{tabular}{@{}lrrr@{}}
\toprule
Teacher arm & Matches & $\Delta E$ & C-sat \\
\midrule
Natural, 0-shot & 81 & 36.89 & 0.650 \\
Natural, 4-shot & 120 & 34.40 & 0.863 \\
Restricted, 0-shot & 104 & 31.94 & 0.867 \\
Restricted, 4-shot & 120 & 33.68 & 0.868 \\
Oracle & 120 & 33.78 & 0.868 \\
\bottomrule
\end{tabular}
\caption{Oracle-assisted direction restatement and one-revision reception,
120 cases on 60 examples. Matches are supplementary resolved set matches;
natural zero-shot has 22 unresolved cases. All receiver endpoints remain in
the comparison, with common starting error 39.83. C-sat uses supplied oracle
directions, not inferred teacher intent.}
\label{tab:teacher_receiver}
\end{table}

Restricted zero-shot reduces receiver error relative to natural zero-shot
by 4.95 [$1.32$, $8.39$]. The natural four-shot mean improvement is 2.49,
but its interval includes zero. Faithful restatement therefore does not
guarantee the best receiving accuracy. Restricted four-shot is close to the
oracle (difference $-0.10$ [$-0.40$, 0.09]), without establishing equivalence.

Restricted zero-shot's advantage over the oracle is concentrated:
11 wins, 104 ties, 5 losses, and median difference zero.
Its 98 literally identical messages tie; six changed same-set messages and
sixteen messages with additions contribute 57.1\% and 42.9\% of its net
advantage. This is an arithmetic decomposition, not a causal effect of
surface wording or extra information. Four of the six same-set messages
reorder directions; two change punctuation or conjunctions. One reordered
message contributes 92.01 of their net 126.04 summed error advantage,
showing concentration in a single case rather than a consistent wording
benefit. The sixteen additions also prevent an
information-matched claim of better restatement.
These controls establish prompt sensitivity in oracle-assisted feedback,
not a general inability to generate geometrically correct guidance unaided.

\section{Discussion}
\label{sec:discussion}
\paragraph{Language as a graded control interface.}
Direction, magnitude, and stopping are separate requirements. Given
$\Delta E_{76}\leq5$, further valid-axis corrections can still overshoot.
A saved-prefix stopping replay retains nine otherwise lost convergence
outcomes and saves 23/576 revisions in the shared-start study
(Appendix~\ref{app:stopping}); this is target-known stopping, not an
available model-only policy. Calibrated wording improves control above the
threshold, but display projection and a sparse phrase ladder limit transfer.

\paragraph{Grounding versus coordinate execution.}
Shared-start results isolate correction from initial name prediction;
isolated numeric probes further separate execution from linguistic magnitude.
LAB with a legend is an instructed coordinate interface, not proof of unaided
perceptual understanding. The observed color control is relevant to creative
tools, but its utility in image editing, where descriptions and transformations
are richer, requires separate evaluation.

\paragraph{Teacher claims require matched evidence.}
Comparative framing, restricted instructions, and demonstrations affect
restatement. A word-level audit is not a semantic judgment, and preserving
directions is not equivalent to maximizing target-error reduction.
Our conclusions concern tested prompt pipelines and receiving behavior,
not a universal following-versus-generating capability gap.

\section{Conclusion and Future Work}
\label{sec:conclusion}
\textsc{ColorRef} evaluates feedback-driven color prediction beyond a
one-shot endpoint. Controlled feedback improves target localization;
matched controls expose encoding sensitivity, direction-dependent
graded steps, and the distinction between teacher restatement and usefulness.
A frozen magnitude-wording policy transfers to fresh controlled targets,
with lower error and evaluation-call cost, while projection drift and numeric
execution tails remain. Future work should study alternative target
distributions, human and multilingual interaction, cross-model transfer,
multi-axis magnitude policies, and downstream creative editing.

\section{Limitations}
\label{sec:limitations}
Our recorded crowd targets need not be uniquely appropriate for their
descriptions. Both the retained corpus and raw English source record one
color per description, without repeated targets for estimating agreement
or acceptable variation. Validity filtering is not target-agreement measurement;
alternative-color judgments or target distributions remain needed,
especially for abstract names. CIELAB is approximately perceptually uniform,
but Euclidean $\Delta E_{76}$ is not a human magnitude calibration or a
strict visibility threshold. Equal coordinate steps have equal norm in
this metric; differences between directions can also reflect generated
native steps and the restricted sRGB gamut, rather than perceptual
nonuniformity alone.

The new controls use one model and prompt family, finite fixed anchors or
twelve-color sampling clusters, English descriptions, and HEX-derived
starting states. Bootstrap intervals do not establish broad population
generalization, generation-randomness uncertainty, or multiplicity-adjusted
significance. Failed calibration and held-out outputs remain in counts;
matched parsed cohorts can introduce
selection bias. No outputs are repaired or filtered by projection cost.

Teacher studies test oracle-assisted restatement, not unaided geometry;
supplementary finite-grammar and assistant review are not independent human
labels. Restriction changes several prompt instructions simultaneously.
Numeric feedback adds exact magnitude and coordinate holds, and calibrated
wording adds a selected magnitude category; neither is information-matched
to bare directions. Target-known stopping and benchmark oracles are not
deployable model-only policies. Evaluation-call savings omit calibration
cost, and fixed-axis transfer does not intentionally correct off-axis drift.
Creative applications, multilingual behavior, HSV magnitude calibration,
and alternative perceptual metrics require further studies.

\begingroup
% Permit flexible bibliography wrapping without changing fonts or margins.
\let\aclthebibliography\thebibliography
\renewcommand{\thebibliography}[1]{%
  \aclthebibliography{#1}\raggedright}
\bibliography{custom}
\endgroup
\appendix
\input{latex/revision-appendix}
\end{document}
